\phantomsection\label{fig:vol01:spring-autumn:states-overview}
\section{空间总览：诸侯怎样围绕中原展开}
\begin{center}
\begin{tikzpicture}[x=.83cm,y=.83cm,font=\small]
  \fill[paper] (0,0) rectangle (12,7);
  \draw[blue!55,line width=1pt] (.2,5.0) .. controls (3.5,4.6) and (5.5,5.2) .. (10.5,4.7);
  \node[text=blue!70!black,above] at (5.3,4.9) {黄河};
  \draw[blue!50,line width=.8pt] (5.0,.5) .. controls (6.3,2.1) and (7.8,2.1) .. (10.4,.9);

  \fill[dynasty!20] (4.7,3.1) -- (5.8,3.2) -- (5.8,4.1) -- (4.9,4.4) -- cycle;
  \node at (5.3,3.75) {周};
  \fill[orange!22] (7.9,4.2) -- (10.0,4.2) -- (10.2,5.7) -- (8.4,5.9) -- cycle;
  \node at (9.1,5.05) {齐};
  \fill[green!20] (5.0,4.6) -- (7.2,4.5) -- (7.4,6.0) -- (5.5,6.2) -- cycle;
  \node at (6.2,5.4) {晋};
  \fill[purple!18] (5.3,1.2) -- (7.7,1.2) -- (8.1,3.4) -- (6.3,3.7) -- cycle;
  \node at (6.7,2.35) {楚};
  \fill[red!15] (1.3,3.1) -- (3.6,3.0) -- (4.0,4.6) -- (2.1,5.0) -- cycle;
  \node at (2.6,3.95) {秦};
  \fill[yellow!25] (6.4,3.7) -- (7.5,3.6) -- (7.7,4.3) -- (6.7,4.5) -- cycle;
  \node[align=center] at (7.0,4.1) {郑\\宋\\鲁};
  \fill[cyan!18] (8.2,.5) -- (10.7,.6) -- (10.6,2.4) -- (8.5,2.7) -- cycle;
  \node at (9.5,1.65) {吴越};

  \draw[-{Stealth[length=2mm]},thick,color=timeline] (9.0,5.0) -- (6.7,4.8);
  \node[above,text=timeline] at (7.8,4.9) {齐的会盟};
  \draw[-{Stealth[length=2mm]},thick,color=dynasty] (6.8,2.8) -- (6.5,4.5);
  \node[right,text=dynasty] at (6.7,3.65) {晋楚竞争};
\end{tikzpicture}

\smallskip
\small\textit{图  春秋主要政治体相对位置示意：用以理解齐、晋、楚、秦及吴越的地缘条件，不表示精确疆界或同时存在的固定国界。}
\end{center}
